\documentclass[aspectratio=169]{beamer}

% ============================================================
% Custom theme replicating the reference deck's look: sans-serif
% body, blue frame titles, hollow/filled triangle bullets, and a
% three-segment colored footline (institution | short title |
% date/page). Built from scratch (no \usetheme{...}) to match
% the reference slides exactly rather than approximate a stock
% Beamer theme.
% ============================================================

\usepackage[T1]{fontenc}
\usepackage{lmodern}
\usepackage[utf8]{inputenc}
\usepackage{amsmath, amssymb}
\usepackage{booktabs}
\usepackage{graphicx}

\setbeamertemplate{navigation symbols}{}

% ---- Colors -------------------------------------------------
\definecolor{slidetitleblue}{RGB}{60,90,190}
\definecolor{footdark}{RGB}{35,45,90}
\definecolor{footmid}{RGB}{70,95,175}

\setbeamercolor{structure}{fg=slidetitleblue}
\setbeamercolor{frametitle}{fg=slidetitleblue,bg=white}
\setbeamercolor{title}{fg=slidetitleblue}
\setbeamercolor{normal text}{fg=black,bg=white}
\setbeamercolor{footline author}{fg=white,bg=footdark}
\setbeamercolor{footline title}{fg=white,bg=footmid}
\setbeamercolor{footline date}{fg=white,bg=footdark}

\setbeamerfont{frametitle}{series=\mdseries,size=\large}
\setbeamerfont{footline author}{size=\scriptsize}
\setbeamerfont{footline title}{size=\scriptsize}
\setbeamerfont{footline date}{size=\scriptsize}

% ---- Bullets: hollow triangle (level 1), filled triangle (level 2) --
\setbeamertemplate{itemize item}{\color{slidetitleblue}$\triangleright$}
\setbeamertemplate{itemize subitem}{\color{slidetitleblue}$\blacktriangleright$}
\setbeamertemplate{itemize subsubitem}{\color{slidetitleblue}$\bullet$}

\setbeamertemplate{frametitle}{%
  \vspace{0.4em}%
  \hspace{0.6em}\usebeamerfont{frametitle}\usebeamercolor[fg]{frametitle}\insertframetitle%
  \vspace{0.3em}%
}

% ---- Three-segment footline ---------------------------------
\newcommand{\shorttitletext}{Sovereign Spread Crises and the Sovereign-Bank Nexus}
\newcommand{\shortauthortext}{David Betsem}

\setbeamertemplate{footline}{%
  \leavevmode%
  \hbox{%
    \begin{beamercolorbox}[wd=.27\paperwidth,ht=2.4ex,dp=1.2ex,left,leftskip=0.6em]{footline author}%
      \usebeamerfont{footline author}\shortauthortext%
    \end{beamercolorbox}%
    \begin{beamercolorbox}[wd=.53\paperwidth,ht=2.4ex,dp=1.2ex,center]{footline title}%
      \usebeamerfont{footline title}\shorttitletext%
    \end{beamercolorbox}%
    \begin{beamercolorbox}[wd=.20\paperwidth,ht=2.4ex,dp=1.2ex,right,rightskip=0.6em]{footline date}%
      \usebeamerfont{footline date}\insertshortdate{} \hspace{0.4em} \insertframenumber{} / \inserttotalframenumber%
    \end{beamercolorbox}%
  }%
}

\setbeamertemplate{navigation symbols}{}
\setbeamertemplate{headline}{}

% ---- Title page (affiliations as superscripts, small disclaimer) ----
\title{\large Costs of Sovereign Debt Crises With and Without Default:\\
Effect of the Sovereign-Bank Nexus}
\author{\normalsize David Betsem A Djam}
\institute{%
  \footnotesize
  HEC Paris, Institut Polytechnique de Paris
}
\date{\normalsize 28\textsuperscript{th} September 2026}

\begin{document}

% ---------------- Title slide ----------------
{
\setbeamertemplate{footline}{}
\begin{frame}[plain]
  \vspace{2cm}
  \vfill
  \begin{center}
    \normalsize Master Thesis Defense\par
  \end{center}

  \begin{center}
    \usebeamercolor[fg]{title}{\Large Costs of Sovereign Debt Crises With and\\
    Without Default: Effect of the Sovereign-Bank Nexus\par}
  \end{center}
  \vspace{0.2cm}
  \begin{center}
    \large David Betsem\par
  \end{center}
  \vspace{0.6em}
  \begin{center}
    \small 28\textsuperscript{th} September 2026\par
  \end{center}
  \vfill
\end{frame}
}

% ---------------- Motivation ----------------
\begin{frame}{Motivation: not all sovereign debt crises end in default}
  \begin{itemize}
    \item \textbf{Sovereign spread crises are more common than defaults}
    \begin{itemize}
      \item Episodes of severe bond-market distress (Mexico 1994/95, Thailand
      1997/98, Russia 2015) occurred without a missed payment or restructuring
      \item Elevated spreads reflect rising default risk even when no default
      is ultimately realized \citep{arellano2008default}
    \end{itemize}
    \item \textbf{The existing literature measures the cost of default}
    \begin{itemize}
      \item Output costs of default are estimated at 1--4 pp of GDP growth
      \item Rarely benchmarked against equally severe episodes that did not
      end in default
    \end{itemize}
    \item \textbf{This paper compares both types of crisis directly}
    \begin{itemize}
      \item Do spread crises without default carry a materially different
      cost than default-linked crises?
      \item Does the pre-crisis sovereign-bank nexus shape that difference?
    \end{itemize}
  \end{itemize}
\end{frame}

% ---------------- This paper ----------------
\begin{frame}{This paper}
  \begin{itemize}
    \item \textbf{Research question:} How do GDP, bank credit, and investment
    respond to sovereign spread crises with and without default, and does the
    pre-crisis sovereign-bank nexus amplify that response?
    \item \textbf{Data:} JP Morgan EMBIG spread database, 1993--2023
    \begin{itemize}
      \item 61 spread crisis episodes, 21 default-linked, 40 non-default
    \end{itemize}
    \item \textbf{Method:} Local projection with Augmented Inverse Probability
    Weighting (AIPW), following \citet{jorda2016time}
    \item \textbf{Identification:} Comparison across crisis type and
    pre-crisis sovereign-bank exposure (median split on bank claims on
    government / assets)
  \end{itemize}
\end{frame}

% ---------------- Preview of results ----------------
\begin{frame}{Preview of results}
  \begin{itemize}
    \item \textbf{GDP and bank credit differ sharply by crisis type}
    \begin{itemize}
      \item GDP falls 1--2 pp following non-default crises vs.\ 6--9 pp
      following default-linked crises
      \item Bank credit turns positive after non-default crises but
      contracts sharply after default-linked crises
    \end{itemize}
    \item \textbf{Investment, FDI, and lending rates respond similarly
    across types}
    \item \textbf{The sovereign-bank nexus matters, but only without default}
    \begin{itemize}
      \item High pre-crisis nexus is associated with a significantly better
      GDP response in non-default crises
      \item No significant high/low-nexus difference in default-linked
      crises, given the smaller sample
    \end{itemize}
  \end{itemize}
\end{frame}

% ---------------- Roadmap ----------------
\begin{frame}{Roadmap}
  \vspace{1cm}
  \begin{enumerate}
    \item \textbf{Data and crisis classification}
    \item Estimation strategy and main results
    \item Role of the sovereign-bank nexus
    \item Conclusion
  \end{enumerate}
\end{frame}

% ---------------- Section divider example ----------------
\begin{frame}{Data sources}
  \begin{itemize}
    \item \textbf{JP Morgan EMBIG Global}
    \begin{itemize}
      \item Sovereign bond spread series used to identify crisis onsets
    \end{itemize}
    \item \textbf{IMF WEO, IMF FSI, World Bank IDS}
    \begin{itemize}
      \item GDP, investment, bank credit, bank claims on government, public
      debt, trade openness
    \end{itemize}
    \item \textbf{\citet{asonuma2016sovereign}}
    \begin{itemize}
      \item Sovereign restructuring database, used to classify episodes as
      default-linked or non-default
    \end{itemize}
    \item \textbf{\citet{laeven2013systemic,laeven2026systemic}}
    \begin{itemize}
      \item Systemic banking crisis dummy
    \end{itemize}
  \end{itemize}
\end{frame}

% ---------------- Specification slide (equation example) ----------------
\begin{frame}{Baseline specification}
  \begin{center}
    \[
    y_{i,t+h} - y_{i,t} = \alpha_i^{S,h} + \beta^{S,h} D_{i,t+1}^{S} + X_{i,t}\beta^{S,h} + \mu_{i,t+h}^{S}
    \]
  \end{center}
  \vspace{0.5em}
  \begin{itemize}
    \item $y_{i,t+h} - y_{i,t}$: cumulative change in the outcome from crisis
    onset to horizon $h$
    \item $D_{i,t+1}^{S}$: onset dummy for crisis type $S \in$ \{non-default,
    default-linked\}
    \item $X_{i,t}$: pre-crisis controls
    \item $\alpha_i^{S,h}$: country fixed effects
    \item Estimated via AIPW to attenuate selection bias, following
    \citet{jorda2016time}
  \end{itemize}
\end{frame}

% ---------------- Table example ----------------
\begin{frame}{GDP response, AIPW estimates}
  \begin{center}
    \footnotesize
    \begin{tabular}{lccccc}
      \toprule
       & $h=0$ & $h=1$ & $h=2$ & $h=3$ & $h=4$ \\
      \midrule
      Non-default & -- & -- & -- & -- & -- \\
      Default-linked & -- & -- & -- & -- & -- \\
      \bottomrule
    \end{tabular}
  \end{center}
  \vspace{0.5em}
  \begin{itemize}
    \item Replace with the estimates and standard errors from
    \texttt{tab:aipw} (Panel A, GDP)
  \end{itemize}
\end{frame}

% ---------------- Conclusion ----------------
\begin{frame}{Conclusion}
  \begin{itemize}
    \item Sovereign spread crises without default carry real but modest
    output costs, far smaller than default-linked crises
    \item The sovereign-bank nexus amplifies the recovery in non-default
    crises, consistent with a bank-lending channel that is only active when
    the sovereign avoids default
    \item No significant nexus effect within default-linked crises, though
    the point estimates are consistently more negative for high-nexus
    countries -- a null result given the limited sample of default-linked
    episodes, not evidence of no effect
  \end{itemize}
\end{frame}

\end{document}
